\documentclass[11pt]{article}
\usepackage[margin=1in]{geometry}
\usepackage[T1]{fontenc}
\usepackage{lmodern}
\usepackage{booktabs}
\usepackage{amsmath}
\usepackage[hidelinks]{hyperref}
\usepackage{xurl}
\title{CSE498 Lab 2: ReVA Video Question Answering}
\author{Ling Shen}
\date{}

\begin{document}
\maketitle
\noindent\textbf{Project repository:} \url{https://github.com/lllingshen/CSE498_lab2}

\section{Setup and implementation}

This lab compares Qwen3-VL-4B-Instruct before and after LoRA fine-tuning,
with VILA1.5-3b as a second baseline. I used the course starter code\footnote{\url{https://github.com/likaiw2/VLM_evaluating_finetuing}}
and the official ReVA annotations.\footnote{\url{https://huggingface.co/datasets/ReVA-Benchmark/ReVA}}
The training subset contains the first 200 training questions. All three
models were evaluated on the same 4,000 test questions from 1,014 videos.

The annotation converter produces a conversation for each training question:
the user turn contains a video reference, the question, and its lettered
options; the assistant turn contains the correct letter inside an
\texttt{<answer>} tag. Test preparation retains each question's QA ID,
video reference, options, answer, and subcategory. The Qwen pipeline evaluates
the base model, trains a LoRA adapter, and evaluates that adapter on the same
test set. VILA uses its own video-question-answering adapter and is evaluated
without fine-tuning.

I completed the core assignment functions in these files:
\begin{itemize}
\setlength{\itemsep}{0pt}
\item \path{scripts/build_qwen_train_data.py}: converts annotations into
video-question conversations and answer targets.
\item \path{reva_eval/data/rsvidqa/prepare_reva_v2_test_set.py}: prepares
test questions with stable QA IDs and video references.
\item \path{scripts/score_reva_predictions.py}: parses Qwen answers and
scores every ground-truth question.
\item \path{vila_eval/reva_v2.py}: loads ReVA questions, builds VILA
prompts, parses answers, and summarizes results.
\end{itemize}
Other local execution changes include
\path{scripts/run_eval_qwen_finetuned.sh},
\path{reva_eval/inference_vllm_origin_number.py}, and
\path{scripts/check_student_setup.py}. These check LoRA adapter loading,
handle video metadata, and support the local environment and data paths.
Single-GPU SDPA training support is implemented in
\path{qwen_finetune/qwenvl/train/train_qwen.py} and
\path{qwen_finetune/scripts/sft_7b.sh}.
Focused regression checks are in \path{tests/test_local_fixes.py}.

The saved training arguments record one epoch, batch size 1, gradient
accumulation of 4, and learning rate $2\times10^{-7}$. LoRA uses rank 8,
alpha 16, and dropout 0 on the attention query, key, value, and output
projections. The training state records 50 optimizer steps and mean training
loss 0.4115. The saved evaluation configuration points to the trained adapter
and the Qwen base model; the run log records adapter loading followed by
successful model and processor loading.

The recorded experiments use a maximum of four frames per video and greedy
decoding. Qwen receives a request for reasoning followed by an answer tag,
with a maximum of 1,024 output tokens. VILA receives a request for only the
option letter, with a maximum of 128 output tokens. These are the settings
used for the saved runs, rather than recommendations for further tuning.
Reproduction commands are in \texttt{README.md}.

Scores were recomputed from the saved raw responses. The saved evaluation
ground truth matches the official test annotations in QA IDs, questions,
options, labels, video references, and subcategories. Each run has all 4,000
unique IDs and no duplicates. The scoring rule is
\[
  \mathrm{Accuracy} =
  \frac{\text{number of correct answers}}{\text{all ground-truth questions}}.
\]
Missing or unparseable answers count as wrong. Completion means that an
answer can be parsed into an option letter; it does not mean the answer is
correct. An explicit final answer takes precedence over distractors mentioned
in the explanation.

\section{Results and interpretation}

\begin{center}
\begin{tabular}{lrrr}
\toprule
Model & Correct / total & Accuracy & Parseable / total \\
\midrule
Qwen base & 2,675 / 4,000 & 66.88\% & 3,971 / 4,000 \\
Qwen LoRA & 2,673 / 4,000 & 66.83\% & 3,964 / 4,000 \\
VILA base & 1,940 / 4,000 & 48.50\% & 4,000 / 4,000 \\
\bottomrule
\end{tabular}
\end{center}

LoRA did not improve overall accuracy in this experiment. It produced two
fewer correct answers, a decrease of 0.05 percentage points. A direct
question-by-question comparison shows that 93 previously wrong answers
became correct, while 95 correct answers became wrong. All models generated
responses for all questions, but 29 base-Qwen responses and 36 LoRA responses
were unparseable and therefore counted as incorrect.

There were some differences by subcategory. In Geometric Relation, Qwen
improved from 270/400 (67.50\%) to 280/400 (70.00\%) after fine-tuning.
Temporal Grounding remained difficult: the base model answered 252/640
(39.38\%) correctly, compared with 250/640 (39.06\%) after fine-tuning.
These local changes did not yield an overall improvement. The training loss
and these scores alone do not establish overfitting or the cause of the
difference. Qwen outperformed VILA under the settings used here, with the
prompt and output-budget differences noted above.

\section{Three qualitative examples}

\paragraph{1. Counting cars: QA-000005 (visual perception).}
The question asks how many cars are visible in the parking lot at 00:00.
The options are A: 40--50, B: 10--20, C: 20--30, and D: 30--40.
The official answer is \textbf{B}; Qwen base predicts \textbf{D}, Qwen LoRA
predicts \textbf{C}, and VILA predicts \textbf{B}.
Both Qwen responses select a count range above the official answer. The
LoRA response changes the range but remains incorrect. This is consistent
with a visual counting difficulty, although the outputs alone do not show
why the models differ. The timestamp matters: the question asks for the
count at the beginning, not a count accumulated across the video.

\newpage
\paragraph{2. Visibility interval: QA-000010 (temporal reasoning).}
The question asks how long a group of people remains visible in the parking
lot. The options are A: 1--5 s, B: 0--4 s, C: 0--5 s, and D: 2--5 s.
The official answer is \textbf{C}; Qwen base and Qwen LoRA both predict
\textbf{C}, while VILA predicts \textbf{D}.
VILA selects a later starting time and a shorter duration than the official
answer. The options differ by only one or two seconds, making this a useful
temporal-reasoning example. With only four sampled frames, the input provides
limited evidence for precise boundaries. Choosing the correct interval does
not by itself demonstrate continuous tracking of the group.

\paragraph{3. Viewpoint and answer mapping: QA-000248 (option parsing).}
The question asks for the camera viewpoint. The options are A: high angle
($>45^\circ$), B: top-down ($90^\circ$), C: high angle ($<45^\circ$), and
D: ground-level ($0^\circ$). The official answer is \textbf{A}; Qwen base
predicts \textbf{C}, Qwen LoRA predicts \textbf{A}, and VILA predicts
\textbf{B}.
The base model's explanation includes ``high angle ($>45$ degrees)'',
but its final tag is \texttt{<answer>C</answer>}. This is a model
answer-mapping inconsistency, not a scorer error: the final letter must be
scored as C even though the explanation describes option A. The LoRA
response aligns its explanation with A. This one corrected example does
not establish an overall benefit from fine-tuning.

\section{Limitations and conclusion}

This was a small, one-epoch training run, and the four-frame budget limits
the available video information. The models also use different prompts and
output limits. The 200 training questions come from 17 videos; 16 of those
videos also occur in the test set, covering 68 test questions. No identical
video/question pair occurs in both this training subset and the test set,
but the evaluation is not entirely on unseen videos. These limitations do
not establish the cause of any score change.

The experiment shows how to convert video annotations, train and load a
LoRA adapter, and compare models using a common test denominator. Qwen
performed better than VILA in this run, while the small LoRA fine-tuning
run left Qwen's overall accuracy essentially unchanged.

\end{document}
